% Created 2026-02-28 Sat 21:52
% Intended LaTeX compiler: pdflatex
\documentclass[8pt]{beamer}

\usepackage[hyperref]{biblatex}
\addbibresource{~/Documents/GLOBAL_ORG/org/roam/bib/master.bib}
\usepackage[utf8]{inputenc}
\usepackage[T1]{fontenc}
\usepackage[italian, english]{babel}
\usefonttheme{professionalfonts}
\usepackage{eurosym}
\DeclareUnicodeCharacter{20AC}{\euro}
\usetheme{Madrid}
\author{Davide Peccioli}
\date{11-20 agosto 2026}
\title{Vacanze ``Ultimo Cena''}
\begin{document}

\maketitle
\begin{frame}[label={sec:org99b306c}]{Introduzione}
\begin{itemize}
\item \href{https://www.airbnb.it/rooms/1243277132509758895?guests=8\&adults=8\&s=67\&unique\_share\_id=aba1fee1-90d0-4b31-bbef-b2b30f9853b1}{Link AirBNB}
\item Periodo: \alert{11-20 agosto 2026} (9 notti)
\item Costo casa: \alert{1805 €}
\item Costo macchina: 128 € (\emph{prezzo stimato solo andata, per ogni macchina}).
\end{itemize}

Per il calcolo preciso dei costi, l'idea è la seguente: il costo a testa per notte deve essere lo stesso, mentre per la macchina si approssima solo nei due grandi viaggi (andata e ritorno), e si divide equamente.

Ho diviso tre scenari possibili, viste le esigenze di Alessia e Roberto.
\end{frame}
\begin{frame}[label={sec:org23fbee0}]{Alessia e Roberto passeranno 5 notti con noi}
Nello specifico:
\begin{itemize}
\item \alert{Alessia e Roberto} passano 5 notti in vacanza (partenza il 16 agosto), e viaggeranno in macchina solo per l'andata:
\begin{align*}
  \text{Casa} &= \frac{1805\ \text{euro}}{(6 \cdot 9\ \text{notti}) + (2\cdot 5\ \text{notti})} \cdot (5\ \text{notti}) = 142 \ \text{euro}\\[2ex]
  \text{Macchina} &= \frac{4\cdot 128\ \text{euro}}{8 + 6} \cdot 1 = 37\ \text{euro}
\end{align*}
Quindi spenderanno \alert{179 euro}.
\item \alert{Tutti gli altri} passano 9 notti in vacanza e fanno sia andata che ritorno:
\begin{align*}
  \text{Casa} &= \frac{1805\ \text{euro}}{(6 \cdot 9\ \text{notti}) + (2\cdot 5\ \text{notti})} \cdot (9\ \text{notti}) = 254\ \text{euro}\\[2ex]
  \text{Macchina} &= \frac{4\cdot 128\ \text{euro}}{8 + 6} \cdot 2 = 74\ \text{euro}
\end{align*}
Quindi spenderanno \alert{328 euro} (+ 21 euro rispetto al preventivo precedente).
\end{itemize}
\end{frame}
\begin{frame}[label={sec:org62086e4}]{Alessia e Roberto passeranno, rispettivamente, 4 e 6 notti con noi.}
Nello specifico:
\begin{itemize}
\item \alert{Alessia} passa 4 notti in vacanza (partenza il 15 agosto), e viaggerà in macchina solo per l'andata:
\begin{align*}
  \text{Casa} &= \frac{1805\ \text{euro}}{(6 \cdot 9\ \text{notti}) + (6\ \text{notti}) + (4\ \text{notti})} \cdot (4\ \text{notti}) =  112\ \text{euro}\\[2ex]
  \text{Macchina} &= \frac{4\cdot 128\ \text{euro}}{8 + 6} \cdot 1 = 37\ \text{euro}
\end{align*}
Quindi spenderà \alert{149 euro}.
\item \alert{Roberto} passa 6 notti in vacanza (partenza il 17 agosto), e viaggerà in macchina solo per l'andata:
\begin{align*}
  \text{Casa} &= \frac{1805\ \text{euro}}{(6 \cdot 9\ \text{notti}) + (6\ \text{notti}) + (4\ \text{notti})} \cdot (6\ \text{notti}) = 170\ \text{euro} \\[2ex]
  \text{Macchina} &= \frac{4\cdot 128\ \text{euro}}{8 + 6} \cdot 1 = 37\ \text{euro}
\end{align*}
Quindi spenderà \alert{207 euro}.
\item \alert{Tutti gli altri} passano 9 notti in vacanza e fanno sia andata che ritorno:
\begin{align*}
  \text{Casa} &= \frac{1805\ \text{euro}}{(6 \cdot 9\ \text{notti}) + (6\ \text{notti}) + (4\ \text{notti})} \cdot (9\ \text{notti}) = 254\ \text{euro}\\[2ex]
  \text{Macchina} &= \frac{4\cdot 128\ \text{euro}}{8 + 6} \cdot 2 = 74\ \text{euro}
\end{align*}
Quindi spenderanno \alert{328 euro} (+ 21 euro rispetto al preventivo precedente).
\end{itemize}
\end{frame}
\begin{frame}[label={sec:orgea7172f}]{Alessia e Roberto passeranno, rispettivamente, 5 e 6 notti con noi.}
Nello specifico:
\begin{itemize}
\item \alert{Alessia} passa 5 notti in vacanza (partenza il 16 agosto), e viaggerà in macchina solo per l'andata:
\begin{align*}
  \text{Casa} &= \frac{1805\ \text{euro}}{(6 \cdot 9\ \text{notti}) + (6\ \text{notti}) + (5\ \text{notti})} \cdot (5\ \text{notti}) = 139 \ \text{euro}\\[2ex]
  \text{Macchina} &= \frac{4\cdot 128\ \text{euro}}{8 + 6} \cdot 1 = 37\ \text{euro}
\end{align*}
Quindi spenderà \alert{176 euro}.
\item \alert{Roberto} passa 6 notti in vacanza (partenza il 17 agosto), e viaggerà in macchina solo per l'andata:
\begin{align*}
  \text{Casa} &= \frac{1805\ \text{euro}}{(6 \cdot 9\ \text{notti}) + (6\ \text{notti}) + (5\ \text{notti})} \cdot (6\ \text{notti}) = 167\ \text{euro} \\[2ex]
  \text{Macchina} &= \frac{4\cdot 128\ \text{euro}}{8 + 6} \cdot 1 = 37\ \text{euro}
\end{align*}
Quindi spenderà \alert{204 euro}.
\item \alert{Tutti gli altri} passano 9 notti in vacanza e fanno sia andata che ritorno:
\begin{align*}
  \text{Casa} &= \frac{1805\ \text{euro}}{(6 \cdot 9\ \text{notti}) + (6\ \text{notti}) + (5\ \text{notti})} \cdot (9\ \text{notti}) = 250\ \text{euro}\\[2ex]
  \text{Macchina} &= \frac{4\cdot 128\ \text{euro}}{8 + 6} \cdot 2 = 74\ \text{euro}
\end{align*}
Quindi spenderanno \alert{324 euro} (+ 17 euro rispetto al preventivo precedente).
\end{itemize}
\end{frame}
\begin{frame}[label={sec:orgfeb33ca}]{Alessia e Roberto passeranno 6 notti con noi}
Nello specifico:
\begin{itemize}
\item \alert{Alessia e Roberto} passano 6 notti in vacanza (partenza il 17 agosto), e viaggeranno in macchina solo per l'andata:
\begin{align*}
  \text{Casa} &= \frac{1805\ \text{euro}}{(6 \cdot 9\ \text{notti}) + (2\cdot 6\ \text{notti})} \cdot (6\ \text{notti}) = 165 \ \text{euro}\\[2ex]
  \text{Macchina} &= \frac{4\cdot 128\ \text{euro}}{8 + 6} \cdot 1 = 37\ \text{euro}
\end{align*}
Quindi spenderanno \alert{202 euro}.
\item \alert{Tutti gli altri} passano 9 notti in vacanza e fanno sia andata che ritorno:
\begin{align*}
  \text{Casa} &= \frac{1805\ \text{euro}}{(6 \cdot 9\ \text{notti}) + (2\cdot 6\ \text{notti})} \cdot (9\ \text{notti}) = 246\ \text{euro}\\[2ex]
  \text{Macchina} &= \frac{4\cdot 128\ \text{euro}}{8 + 6} \cdot 2 = 74\ \text{euro}
\end{align*}
Quindi spenderanno \alert{320 euro} (+ 13 euro rispetto al preventivo precedente).
\end{itemize}
\end{frame}
\end{document}
